% Paper 6 (GRPO-Registry) schema section
\section{Registry Schema}
\label{sec:p6-schema}

The schema (\texttt{registry/schema.json}) defines two record types under one
\texttt{oneOf}. Both carry a \texttt{schema_version} for forward migration
(Section~\ref{sec:p6-governance}). A single field convention governs
everything: \textbf{JSON \texttt{null} means \emph{unreported}}, while explicit
\texttt{false} / \texttt{0.0} / \texttt{"none"} means
\emph{reported-as-absent}. A stack that says ``we use no KL term'' is making a
checkable claim; a stack that says nothing is not. The auditor badge
(Section~\ref{sec:p6-query}) scores exactly this distinction: reporting
coverage, not configuration virtue.

\subsection{Stack Records: the Seven-Item MIN-REPORT Block}
\label{sec:p6-schema-stack}
A \texttt{stack} record describes one deployed GRPO-family training stack. Its
core is the \texttt{min_report} object, one sub-object per item of the
MIN-REPORT-RL minimum reportable stack:

\begin{enumerate}[leftmargin=1.8em, itemsep=0.1em]
  \item \textbf{loss_form}: importance-ratio level (token/sequence), clip
        bounds ($\eps_{\text{low}}, \eps_{\text{high}}$), length
        normalization, advantage normalization (std/none), token mask. These
        are the fields the variant papers fight over.
  \item \textbf{reference_kl}: whether a reference policy exists, the KL
        coefficient $\beta$, and the estimator/surrogate used.
  \item \textbf{sampler_backend}: backend identity, numeric precision, and
        sampler parameters (temperature, top-$p$)---the axis along which our
        backend-swap audit produced the $84.4\%$-vs-$5.0\%$ swing.
  \item \textbf{telemetry}: whether the stack emits a \emph{per-step} ZVF and
        gradient-utilization ($\mathrm{GU}=1-\mathrm{ZVF}$) trajectory, and
        from what source. Endpoint summaries are not enough: our 368-run audit
        found ZVF is U-shaped in accuracy, so a single mean aliases mastery
        with incapacity (see the companion ZVF pillar paper).
  \item \textbf{group_size_schedule}: initial $G$, fixed vs.\ adaptive, and
        the adaptation rule if any.
  \item \textbf{heldout_split}: whether evaluation prompts are disjoint from
        the reward environment, and how the split was built.
  \item \textbf{decontamination}: decontamination performed, and whether a
        parser-robustness probe was run against the reward extractor.
\end{enumerate}

Around the \texttt{min_report} block sit provenance
(\texttt{source_artifacts}, W\&B project, date), identity
(\texttt{label_claimed}, framework name/version, and an
\texttt{openness} enum: \texttt{open}/\texttt{managed}/\texttt{closed}), an
optional \texttt{outcomes} summary (informational only---never part of the
badge), and \texttt{variant_deltas_applied}, described next.

\subsection{Variant-Delta Records}
\label{sec:p6-schema-delta}
A \texttt{variant_delta} record decomposes a named variant into components,
each mapping a schema field to a change relative to base GRPO. We ship three,
each verified against its source paper rather than against folklore:

\begin{itemize}[leftmargin=1.5em, itemsep=0.15em]
  \item \textbf{DAPO}~\citep{yu2025dapo}: five components---%
        \texttt{clip_higher} (asymmetric/decoupled clip,
        $\eps_{\text{low}}{=}0.2$, $\eps_{\text{high}}{=}0.28$),
        \texttt{dynamic_sampling} (filter zero-variance groups and resample
        until the batch is full), \texttt{token_level_loss} (token-level
        policy-gradient aggregation), \texttt{overlong_reward_shaping} (soft
        length-aware penalty), and \texttt{kl_removed} (the KL term is
        excluded from the objective).
  \item \textbf{Dr.~GRPO}~\citep{tong2025drgrpo}: two removals---the
        $1/|o_i|$ response-length normalization in the loss, and the division
        by the within-group reward standard deviation in the advantage.
  \item \textbf{GSPO}~\citep{qwen2025gspo}: the importance ratio is defined on
        (length-normalized) sequence likelihood rather than per token, with
        clipping and optimization applied at the sequence level.
\end{itemize}

\subsection{Where Label Flips Become Machine-Readable}
\label{sec:p6-schema-status}
The joint between the two record types is the stack-side
\texttt{variant_deltas_applied} array: for each component of each claimed
delta, the stack declares a status in
$\{\texttt{implemented}, \texttt{surrogate}, \texttt{absent},
\texttt{unknown}\}$. The shipped ``DAPO-on-managed-runtime'' entry, for
example, records \texttt{clip_higher: surrogate},
\texttt{dynamic_sampling: absent}, and \texttt{token_level_loss: unknown}
(\texttt{managed_by_tinker}):

\begin{lstlisting}[basicstyle=\ttfamily\scriptsize, frame=single,
  caption={Excerpt from \texttt{entries/tinker_dapo_qwen3.5-4b_gsm8k.json}.},
  label={lst:p6-dapo-entry}]
"label_claimed": "dapo",
"variant_deltas_applied": [
  {"delta_id": "delta_dapo", "component": "clip_higher",
   "status": "surrogate",
   "note": "asymmetric clip set through the user-facing config;
            underlying managed loss form unverifiable"},
  {"delta_id": "delta_dapo", "component": "dynamic_sampling",
   "status": "absent",
   "note": "closed sampler exposes no group filter-and-resample hook"},
  {"delta_id": "delta_dapo", "component": "token_level_loss",
   "status": "unknown", "note": "managed_by_tinker"} ]
\end{lstlisting}

A reader---or a script---can now see at a glance that this stack's ``DAPO'' is
one component out of five, without reading a paper appendix or reverse
engineering a launcher.
